\documentclass[11pt]{ctexart}
\usepackage[a4paper,margin=1in]{geometry}
\usepackage{longtable,booktabs}
\title{Geant4-Agent 回归测试报告（中文）}
\author{自动化评测流程}
\date{2026-03-08\_121732}
\begin{document}
\maketitle
\section{统一 Pass/Fail 总表}
\begin{longtable}{p{0.06\linewidth}p{0.08\linewidth}p{0.07\linewidth}p{0.08\linewidth}p{0.07\linewidth}p{0.07\linewidth}p{0.07\linewidth}p{0.07\linewidth}p{0.09\linewidth}p{0.06\linewidth}p{0.21\linewidth}}
\toprule
ID & 领域 & 风格 & 参数完整性 & 期望初始 & 实际初始 & 期望最终 & 实际最终 & 缺参(初->终) & 结果 & 失败原因 \\
\midrule
MTM01 & medical & fuzzy & missing & 失败 & 失败 & 通过 & 通过 & 12->0 & 通过 & - \\
MTM02 & aerospace & precise & missing & 失败 & 失败 & 通过 & 通过 & 3->0 & 通过 & - \\
MTM03 & industrial\_ndt & precise & missing & 失败 & 失败 & 通过 & 通过 & 12->0 & 通过 & - \\
MTM04 & physics & precise & missing & 失败 & 失败 & 通过 & 通过 & 8->0 & 通过 & - \\
MTM05 & nuclear\_engineering & fuzzy & missing & 失败 & 失败 & 通过 & 通过 & 9->0 & 通过 & - \\
MTM06 & security\_screening & fuzzy & missing & 失败 & 失败 & 通过 & 通过 & 11->0 & 通过 & - \\
MTM07 & semiconductor & precise & missing & 失败 & 失败 & 通过 & 通过 & 3->0 & 通过 & - \\
MTM08 & education & fuzzy & missing & 失败 & 失败 & 通过 & 通过 & 11->0 & 通过 & - \\
MTM09 & environmental\_radiation & fuzzy & missing & 失败 & 失败 & 通过 & 通过 & 3->0 & 通过 & - \\
MTM10 & high\_energy\_physics & precise & missing & 失败 & 失败 & 通过 & 通过 & 10->0 & 通过 & - \\
MTM11 & materials\_update & precise & full & 通过 & 通过 & 通过 & 通过 & 0->0 & 通过 & - \\
MTM12 & overwrite\_reject & precise & full & 通过 & 通过 & 通过 & 通过 & 0->0 & 通过 & - \\
MTM13 & medical\_detector & fuzzy & missing & 失败 & 失败 & 通过 & 失败 & 13->1 & 失败 & final\_complete\_mismatch \\
MTM14 & industrial\_panel & precise & missing & 失败 & 失败 & 通过 & 通过 & 15->0 & 失败 & expected\_final\_mismatch:materials.selected\_materials.0 \\
MTM15 & shielding\_stack & precise & missing & 失败 & 失败 & 通过 & 通过 & 14->0 & 通过 & - \\
MTM16 & nested\_target & precise & missing & 失败 & 失败 & 通过 & 通过 & 13->0 & 通过 & - \\
MTM17 & coaxial\_shield & fuzzy & missing & 失败 & 失败 & 通过 & 通过 & 14->0 & 通过 & - \\
MTM18 & boolean\_cavity & precise & missing & 失败 & 失败 & 通过 & 失败 & 17->1 & 失败 & final\_complete\_mismatch, expected\_final\_mismatch:materials.selected\_materials.0 \\
MTM19 & ring\_complete & precise & full & 通过 & 失败 & 通过 & 失败 & 1->1 & 失败 & initial\_complete\_mismatch, final\_complete\_mismatch \\
MTM20 & grid\_complete & precise & full & 通过 & 失败 & 通过 & 失败 & 1->1 & 失败 & initial\_complete\_mismatch, final\_complete\_mismatch, expected\_final\_mismatch:geometry.structure, expected\_final\_mismatch:materials.selected\_materials.0 \\
MTM21 & layered\_detector & fuzzy & missing & 失败 & 失败 & 通过 & 通过 & 14->0 & 通过 & - \\
MTM22 & embedded\_sensor & precise & full & 通过 & 通过 & 通过 & 通过 & 0->0 & 通过 & - \\
MTM23 & layered\_shield\_complete & precise & full & 通过 & 通过 & 通过 & 通过 & 0->0 & 通过 & - \\
MTM24 & boolean\_complete & precise & full & 通过 & 失败 & 通过 & 失败 & 1->1 & 失败 & initial\_complete\_mismatch, final\_complete\_mismatch, expected\_final\_mismatch:materials.selected\_materials.0, expected\_final\_mismatch:source.type \\
\bottomrule
\end{longtable}
\section{基准对比}
\begin{itemize}
\item pass\_rate=0.750
\item initial\_alignment\_rate=0.875
\item final\_alignment\_rate=0.792
\item avg\_turns\_per\_case=2.458
\item missing\_reduction\_rate=1.000
\item confirm\_apply\_success\_rate=1.000
\item rollback\_after\_confirm\_rate=0.000
\item internal\_leak\_turn\_rate=0.051
\item repeated\_assistant\_case\_rate=0.000
\item stale\_confirm\_case\_rate=0.000
\item delta\_pass\_rate=-0.200
\item regressions=无
\item improvements=无
\end{itemize}
\appendix
\section{附录：完整对话与内部处理链（自然语言）}
\subsection{案例 MTM01（medical）}
\textbf{对话：}
\begin{itemize}
\item 用户：我先要一个用于医学成像的初步辐照场景，细节还没定，你先告诉我还缺什么。
\item 系统：目前我们已经确定了使用盒体作为几何结构，接下来还需要补充一些具体参数：盒体在x、y、z三个方向的尺寸，以及为体积指定材料，并选择源的类型。
\item 用户：采用 300 mm x 300 mm x 300 mm 的水箱，材料 ?；源用束流 gamma，能量 6 MeV，位置 (0,0,-400) mm，方向 +z；物理列表 QBBC；输出 root。
\item 系统：我注意到这可能会覆盖已经确认的设置。具体来说：几何结构类型将从“未知”改为“单盒体”，同时盒体 x 方向尺寸将被设置为 300.0。

如果你想应用这些修改，请直接回复“确认”；如果想保留原来的值，请回复“保持原值”。
\item 用户：确认
\item 系统：配置已完成！目前已经确认的设置有：几何结构类型、盒体在 x、y、z 三个方向的尺寸。
\end{itemize}
\textbf{内部链路（自然语言）：}
\begin{itemize}
\item 第 1 轮：动作=summarize\_progress，阶段=geometry，闭环状态=False。
\item 缺失字段：geometry.params.module\_x、geometry.params.module\_y、geometry.params.module\_z、materials.volume\_material\_map、source.type、source.particle、source.energy、source.position、source.direction、physics.physics\_list、output.format、output.path。
\item LLM 使用=True；回退=无；后端=llm\_semantic\_frame+runtime\_semantic。
\item 第 2 轮：动作=confirm\_overwrite，阶段=geometry，闭环状态=False。
\item 缺失字段：geometry.params.module\_x、geometry.params.module\_y、geometry.params.module\_z。
\item LLM 使用=True；回退=无；后端=llm\_slot\_frame+runtime\_semantic。
\item 第 3 轮：动作=finalize，阶段=materials，闭环状态=True。
\item 缺失字段：无。
\item LLM 使用=True；回退=无；后端=llm\_slot\_frame+runtime\_semantic+confirmed\_pending\_overwrite。
\end{itemize}
\subsection{案例 MTM02（aerospace）}
\textbf{对话：}
\begin{itemize}
\item 用户：I need a first-pass proton irradiation setup for a silicon device; keep it simple for now.
\item 系统：Here's a summary of your setup so far and what's still needed.

**Completed:**
*   **Geometry:** A box structure with its X and Y dimensions defined.
*   **Materials:** Materials have been selected and assigned to your geometry volume.
*   **Source:** A proton source with its energy is configured.

**Still Required:**
*   **Physics:** You need to select a physics list.
*   **Output:** You need to specify the output format and the file path to save the data.
\item 用户：Use a silicon box 50 mm x 50 mm x 5 mm, beam source proton 100 MeV from (0,0,-100) mm along +z, physics FTFP\_BERT, output json.
\item 系统：You've requested a change that would overwrite an already confirmed value.

Pending changes:
*   Box size along X: 1000.0 mm → 50.0 mm
*   Box size along Y: 1000.0 mm → 50.0 mm

Please reply **'confirm'** to apply these changes, or **'keep existing'** to keep the current values.
\item 用户：confirm
\item 系统：Your configuration is complete. The following items are confirmed: box size along Z, box size along Y, box size along X, and the output path.
\end{itemize}
\textbf{内部链路（自然语言）：}
\begin{itemize}
\item 第 1 轮：动作=summarize\_progress，阶段=materials，闭环状态=False。
\item 缺失字段：physics.physics\_list、output.format、output.path。
\item LLM 使用=True；回退=无；后端=llm\_slot\_frame+runtime\_semantic。
\item 第 2 轮：动作=confirm\_overwrite，阶段=source，闭环状态=False。
\item 缺失字段：无。
\item LLM 使用=True；回退=无；后端=llm\_slot\_frame+runtime\_semantic。
\item 第 3 轮：动作=finalize，阶段=physics，闭环状态=True。
\item 缺失字段：无。
\item LLM 使用=True；回退=无；后端=llm\_semantic\_frame+runtime\_semantic+confirmed\_pending\_overwrite。
\end{itemize}
\subsection{案例 MTM03（industrial\_ndt）}
\textbf{对话：}
\begin{itemize}
\item 用户：我要做一个 gamma 穿透铁靶的小测试，但我还没把完整参数说完。
\item 系统：目前我们已确认了材料（铁）和粒子类型（伽马射线）。接下来还需要你提供几何结构的具体类型和尺寸，以及体积与材料的绑定方式。请继续补充这些信息，我们就能一步步完成配置了。
\item 用户：请设为半径 40 mm、半长 80 mm 的铁圆柱，材料 ?；点源 gamma，能量 2 MeV，位置 (0,0,-200) mm，方向 +z；物理列表 Shielding；输出 csv。
\item 系统：配置已完成。当前已确认的设置包括：体积与材料绑定、输出路径、输出格式和物理列表。
\end{itemize}
\textbf{内部链路（自然语言）：}
\begin{itemize}
\item 第 1 轮：动作=summarize\_progress，阶段=geometry，闭环状态=False。
\item 缺失字段：geometry.structure、geometry.params.module\_x、geometry.params.module\_y、geometry.params.module\_z、materials.volume\_material\_map、source.type、source.energy、source.position、source.direction、physics.physics\_list、output.format、output.path。
\item LLM 使用=True；回退=无；后端=llm\_slot\_frame+runtime\_semantic。
\item 第 2 轮：动作=finalize，阶段=materials，闭环状态=True。
\item 缺失字段：无。
\item LLM 使用=True；回退=无；后端=llm\_slot\_frame+runtime\_semantic。
\end{itemize}
\subsection{案例 MTM04（physics）}
\textbf{对话：}
\begin{itemize}
\item 用户：Gamma point source 1 MeV at (0,0,-100) toward +z; please build the rest.
\item 系统：Geometry: still needs geometry structure, box size along x. Materials: still needs volume-to-material binding. Source: updated source type, particle type, source energy. Physics: still needs physics list. Output: still needs output format, output path.
\item 用户：Set copper box 1 m x 1 m x 1 m, physics FTFP\_BERT, output hdf5.
\item 系统：Your configuration is complete. The following items are confirmed: volume-to-material binding, output path, output format, and physics list.
\end{itemize}
\textbf{内部链路（自然语言）：}
\begin{itemize}
\item 第 1 轮：动作=summarize\_progress，阶段=geometry，闭环状态=False。
\item 缺失字段：geometry.structure、geometry.params.module\_x、geometry.params.module\_y、geometry.params.module\_z、materials.volume\_material\_map、physics.physics\_list、output.format、output.path。
\item LLM 使用=True；回退=无；后端=llm\_slot\_frame+runtime\_semantic。
\item 第 2 轮：动作=finalize，阶段=materials，闭环状态=True。
\item 缺失字段：无。
\item LLM 使用=True；回退=无；后端=llm\_slot\_frame+runtime\_semantic。
\end{itemize}
\subsection{案例 MTM05（nuclear\_engineering）}
\textbf{对话：}
\begin{itemize}
\item 用户：我想做一个中子屏蔽的概念验证，现在先给一个缺参开头。
\item 系统：目前我们已经确定了几个关键部分：几何结构类型、材料、粒子类型和物理列表都已同步。接下来还需要补充盒体的尺寸（x、y、z方向）、源的类型和能量，以及输出相关的设置。
\item 用户：使用 200 mm x 200 mm x 200 mm 的铅块，材料 ?；点源 neutron，能量 14 MeV，位置 (0,0,-50) mm，方向 +z；物理列表 QGSP\_BERT\_HP；输出 xml。
\item 系统：我检测到这会覆盖已经确认的值。待修改内容：几何结构类型：unknown -> single\_box；已选材料：['G4\_Cu'] -> ['G4\_Pb']。如果要应用这些修改，请直接回复“确认”；如果要保留原值，请直接说“保持原值”。
\item 用户：确认
\item 系统：配置已完成。当前已确认的设置包括：体积与材料绑定、盒体 z 方向尺寸、盒体 y 方向尺寸、盒体 x 方向尺寸。
\end{itemize}
\textbf{内部链路（自然语言）：}
\begin{itemize}
\item 第 1 轮：动作=summarize\_progress，阶段=geometry，闭环状态=False。
\item 缺失字段：geometry.params.module\_x、geometry.params.module\_y、geometry.params.module\_z、source.type、source.energy、source.position、source.direction、output.format、output.path。
\item LLM 使用=True；回退=无；后端=llm\_slot\_frame+runtime\_semantic。
\item 第 2 轮：动作=confirm\_overwrite，阶段=geometry，闭环状态=False。
\item 缺失字段：geometry.params.module\_x、geometry.params.module\_y、geometry.params.module\_z。
\item LLM 使用=True；回退=无；后端=llm\_slot\_frame+runtime\_semantic。
\item 第 3 轮：动作=finalize，阶段=materials，闭环状态=True。
\item 缺失字段：无。
\item LLM 使用=True；回退=无；后端=llm\_semantic\_frame+runtime\_semantic+confirmed\_pending\_overwrite。
\end{itemize}
\subsection{案例 MTM06（security\_screening）}
\textbf{对话：}
\begin{itemize}
\item 用户：Need a quick cargo-screening toy setup with forward photons, but I have not fixed the target shape, material, source energy, or output format yet.
\item 系统：Geometry: updated geometry structure; still needs box size along x, box size along y. Materials: still needs volume-to-material binding. Source: updated particle type; still needs source type, source energy. Physics: still needs physics list. Output: still needs output format, output path.
\item 用户：Use steel cylinder target radius 60 mm half-length 120 mm, beam gamma 3 MeV from (0,0,-500) mm to +z, physics Shielding, output root.
\item 系统：This would overwrite an already confirmed value. Pending change: geometry structure: unknown -> single\_tubs; box size along x: None -> 1000.0. Reply 'confirm' to apply it, or say 'keep existing' to keep the current value.
\item 用户：confirm
\item 系统：Configuration complete. The following items are confirmed: volume-to-material binding, half length, radius, and box size along z.
\end{itemize}
\textbf{内部链路（自然语言）：}
\begin{itemize}
\item 第 1 轮：动作=summarize\_progress，阶段=geometry，闭环状态=False。
\item 缺失字段：geometry.params.module\_x、geometry.params.module\_y、geometry.params.module\_z、materials.volume\_material\_map、source.type、source.energy、source.position、source.direction、physics.physics\_list、output.format、output.path。
\item LLM 使用=True；回退=无；后端=llm\_slot\_frame+runtime\_semantic。
\item 第 2 轮：动作=confirm\_overwrite，阶段=geometry，闭环状态=False。
\item 缺失字段：geometry.params.module\_x、geometry.params.module\_y、geometry.params.module\_z。
\item LLM 使用=True；回退=无；后端=llm\_slot\_frame+runtime\_semantic。
\item 第 3 轮：动作=finalize，阶段=materials，闭环状态=True。
\item 缺失字段：无。
\item LLM 使用=True；回退=无；后端=llm\_semantic\_frame+runtime\_semantic+confirmed\_pending\_overwrite。
\end{itemize}
\subsection{案例 MTM07（semiconductor）}
\textbf{对话：}
\begin{itemize}
\item 用户：我想搭一个硅薄片的中子照射测试，但先只给一个简述。
\item 系统：好的，我们已同步了以下配置：
*   **几何**：盒体的类型与X、Y方向尺寸。
*   **材料**：已选定材料并完成绑定。
*   **源**：粒子类型、能量和位置。

接下来，还需要您补充以下信息才能完成配置：
*   物理列表
*   输出格式
*   输出路径
\item 用户：使用 20 mm x 20 mm x 1 mm 的硅片，材料 ?；各向同性 neutron，能量 2 MeV，位置 (0,0,0) mm；物理列表 QGSP\_BERT\_HP；输出 json。
\item 系统：我注意到新的尺寸设置会覆盖已确认的数值。具体变动如下：
- 盒体 x 方向尺寸：1000.0 → 20.0
- 盒体 y 方向尺寸：1000.0 → 20.0

请确认您的选择：直接回复“确认”以应用修改，或回复“保持原值”以保留现有设置。
\item 用户：确认
\item 系统：配置已完成。当前已确认的参数包括：源位置、源能量、源类型以及盒体的 z 方向尺寸。
\end{itemize}
\textbf{内部链路（自然语言）：}
\begin{itemize}
\item 第 1 轮：动作=summarize\_progress，阶段=materials，闭环状态=False。
\item 缺失字段：physics.physics\_list、output.format、output.path。
\item LLM 使用=True；回退=无；后端=llm\_slot\_frame+runtime\_semantic。
\item 第 2 轮：动作=confirm\_overwrite，阶段=source，闭环状态=False。
\item 缺失字段：无。
\item LLM 使用=True；回退=无；后端=llm\_slot\_frame+runtime\_semantic。
\item 第 3 轮：动作=finalize，阶段=physics，闭环状态=True。
\item 缺失字段：无。
\item LLM 使用=True；回退=无；后端=llm\_semantic\_frame+runtime\_semantic+confirmed\_pending\_overwrite。
\end{itemize}
\subsection{案例 MTM08（education）}
\textbf{对话：}
\begin{itemize}
\item 用户：For classroom demo, I want gamma attenuation in copper; you can ask only essential missing parts.
\item 系统：Here's a summary of your setup so far:

**Completed:**
*   **Materials:** Copper is selected.
*   **Source:** Particle type is set to gamma.
*   **Physics:** A physics list has been chosen.

**Still needed:**
*   **Geometry:** Structure type and box dimensions (X, Y, Z).
*   **Materials:** Assignment of the material to your geometry volume.
*   **Source:** Type, energy, and position.
*   **Output:** Format and file path.
\item 用户：Okay, copper box 100 mm x 100 mm x 100 mm, point gamma 1 MeV at (0,0,-200) mm, direction +z, physics FTFP\_BERT, output csv.
\item 系统：This would change a setting you've already confirmed.

**Change:** Physics list from "Shielding" to "FTFP\_BERT".

Please reply **'confirm'** to apply this change, or **'keep existing'** to keep the current "Shielding" list.
\item 用户：confirm
\item 系统：Your configuration is complete. The following items are confirmed: physics list, output path, output format, and source direction.
\end{itemize}
\textbf{内部链路（自然语言）：}
\begin{itemize}
\item 第 1 轮：动作=summarize\_progress，阶段=geometry，闭环状态=False。
\item 缺失字段：geometry.structure、geometry.params.module\_x、geometry.params.module\_y、geometry.params.module\_z、materials.volume\_material\_map、source.type、source.energy、source.position、source.direction、output.format、output.path。
\item LLM 使用=True；回退=无；后端=llm\_slot\_frame+runtime\_semantic。
\item 第 2 轮：动作=confirm\_overwrite，阶段=materials，闭环状态=False。
\item 缺失字段：无。
\item LLM 使用=True；回退=无；后端=llm\_slot\_frame+runtime\_semantic。
\item 第 3 轮：动作=finalize，阶段=source，闭环状态=True。
\item 缺失字段：无。
\item LLM 使用=True；回退=无；后端=llm\_semantic\_frame+runtime\_semantic+confirmed\_pending\_overwrite。
\end{itemize}
\subsection{案例 MTM09（environmental\_radiation）}
\textbf{对话：}
\begin{itemize}
\item 用户：我想先搭一个空气体积中的低能 gamma 传播场景，细节稍后补。
\item 系统：好的，目前配置进度如下：

**已设置完成：**
*   **几何**：已确定使用盒体，并设置了X、Y方向的尺寸。
*   **材料**：已选定材料，并完成了与体积的绑定。
*   **源**：已设定粒子类型、能量和位置。

**仍需补充以下信息才能生成脚本：**
*   **物理**：请选择物理列表。
*   **输出**：请指定输出格式和文件保存路径。

请提供这些信息以继续。
\item 用户：设为 500 mm x 500 mm x 500 mm 的空气盒，材料 ??；点源 gamma，能量 0.5 MeV，位置 (0,0,-300) mm，方向 +z；物理列表 QBBC；输出 root。
\item 系统：我注意到这会覆盖已确认的设置。具体修改如下：

- **盒体 x 方向尺寸**：1000.0 → 500.0
- **盒体 y 方向尺寸**：1000.0 → 500.0

请确认您的选择：
- 如果应用这些修改，请直接回复“确认”。
- 如果保留原值，请直接回复“保持原值”。
\item 用户：确认
\item 系统：好的，配置已完成。当前已确认的设置包括：源位置、源能量、盒体在Z、Y、X三个方向的尺寸。
\end{itemize}
\textbf{内部链路（自然语言）：}
\begin{itemize}
\item 第 1 轮：动作=summarize\_progress，阶段=materials，闭环状态=False。
\item 缺失字段：physics.physics\_list、output.format、output.path。
\item LLM 使用=True；回退=无；后端=llm\_slot\_frame+runtime\_semantic。
\item 第 2 轮：动作=confirm\_overwrite，阶段=source，闭环状态=False。
\item 缺失字段：无。
\item LLM 使用=True；回退=无；后端=llm\_slot\_frame+runtime\_semantic。
\item 第 3 轮：动作=finalize，阶段=physics，闭环状态=True。
\item 缺失字段：无。
\item LLM 使用=True；回退=无；后端=llm\_semantic\_frame+runtime\_semantic+confirmed\_pending\_overwrite。
\end{itemize}
\subsection{案例 MTM10（high\_energy\_physics）}
\textbf{对话：}
\begin{itemize}
\item 用户：Set geometry as a sphere target; I have not decided radius or source details yet.
\item 系统：Here's where we stand:

**Completed:**
*   **Geometry:** You've set up a sphere target and provided its dimensions (box sizes for the world volume).

**Still Needed:**
*   **Geometry:** The sphere's radius.
*   **Materials:** Which material to assign to the sphere.
*   **Source:** The source type, particle type, and its energy.
*   **Physics:** The physics list to use.
*   **Output:** The output format and file path.

What would you like to provide next?
\item 用户：Sphere target radius 80 mm, lead material, electron point source 5 MeV at (0,0,-200) mm toward +z, physics FTFP\_BERT, output hdf5.
\item 系统：Your configuration is complete. All required settings, including the volume-to-material binding, output path, output format, and physics list, have been confirmed.
\end{itemize}
\textbf{内部链路（自然语言）：}
\begin{itemize}
\item 第 1 轮：动作=summarize\_progress，阶段=geometry，闭环状态=False。
\item 缺失字段：geometry.params.child\_rmax、materials.volume\_material\_map、source.type、source.particle、source.energy、source.position、source.direction、physics.physics\_list、output.format、output.path。
\item LLM 使用=True；回退=无；后端=llm\_slot\_frame+runtime\_semantic。
\item 第 2 轮：动作=finalize，阶段=materials，闭环状态=True。
\item 缺失字段：无。
\item LLM 使用=True；回退=无；后端=llm\_slot\_frame+runtime\_semantic。
\end{itemize}
\subsection{案例 MTM11（materials\_update）}
\textbf{对话：}
\begin{itemize}
\item 用户：Create a complete setup: 100 mm x 100 mm x 100 mm copper box, point gamma 1 MeV at (0,0,-200) mm toward +z, physics QBBC, output json.
\item 系统：Your configuration is complete. The following items have been confirmed: volume-to-material binding, output path, and output format.
\item 用户：Change the target material to aluminum and change the output format to root.
\item 系统：You've requested changes that would overwrite already confirmed values:

*   **Selected materials:** Copper → Aluminum
*   **Output format:** JSON → ROOT

Please reply **'confirm'** to apply these changes, or **'keep existing'** to keep the current values.
\item 用户：confirm
\item 系统：Your configuration is complete. The following items have been confirmed: selected materials, output format, output path, and volume-to-material binding.
\end{itemize}
\textbf{内部链路（自然语言）：}
\begin{itemize}
\item 第 1 轮：动作=finalize，阶段=materials，闭环状态=True。
\item 缺失字段：无。
\item LLM 使用=True；回退=无；后端=llm\_slot\_frame+runtime\_semantic。
\item 第 2 轮：动作=confirm\_overwrite，阶段=source，闭环状态=False。
\item 缺失字段：无。
\item LLM 使用=True；回退=无；后端=llm\_slot\_frame+runtime\_semantic。
\item 第 3 轮：动作=finalize，阶段=physics，闭环状态=True。
\item 缺失字段：无。
\item LLM 使用=True；回退=无；后端=llm\_semantic\_frame+runtime\_semantic+confirmed\_pending\_overwrite。
\end{itemize}
\subsection{案例 MTM12（overwrite\_reject）}
\textbf{对话：}
\begin{itemize}
\item 用户：请创建一个完整配置：100 mm x 100 mm x 100 mm 的铜盒，材料 ?；点源 gamma，能量 1 MeV，位置 (0,0,-200) mm，方向 +z；物理列表 QBBC；输出 json。
\item 系统：配置已完成。当前已确认的设置包括：体积与材料绑定、输出路径和输出格式。
\item 用户：把材料改成 ?，并把输出格式改成 root。
\item 系统：我注意到这次修改会覆盖已确认的设置。具体来说：输出格式将从 json 改为 root。

请确认你的选择：直接回复“确认”来应用修改，或回复“保持原值”以保留原来的设置。
\item 用户：保持原值
\item 系统：好的，输出格式仍将保持为 json，没有更改。
\end{itemize}
\textbf{内部链路（自然语言）：}
\begin{itemize}
\item 第 1 轮：动作=finalize，阶段=materials，闭环状态=True。
\item 缺失字段：无。
\item LLM 使用=True；回退=无；后端=llm\_slot\_frame+runtime\_semantic。
\item 第 2 轮：动作=confirm\_overwrite，阶段=source，闭环状态=False。
\item 缺失字段：geometry.ask.boolean.solid\_a\_size、geometry.ask.boolean.solid\_b\_size。
\item LLM 使用=True；回退=无；后端=llm\_slot\_frame+runtime\_semantic。
\item 第 3 轮：动作=reject\_overwrite，阶段=physics，闭环状态=True。
\item 缺失字段：无。
\item LLM 使用=True；回退=无；后端=llm\_semantic\_frame+runtime\_semantic+rejected\_pending\_overwrite。
\end{itemize}
\subsection{案例 MTM13（medical\_detector）}
\textbf{对话：}
\begin{itemize}
\item 用户：I want a ring of detector modules around a small source region, but I have not fixed the ring radius or module size yet.
\item 系统：Your geometry structure is now set for a ring of detector modules around a small source region.

Here’s what’s been completed and what is still needed next:

**Completed:**
*   **Geometry:** The structure (a ring of modules) is confirmed. The overall bounding box size is also set.

**Still Needed:**
*   **Geometry:** Ring module size, ring module count, and ring radius.
*   **Materials:** Which material to assign to each volume.
*   **Source:** Source type, particle type, and its energy and position.
*   **Physics:** Which physics list to use.
*   **Output:** The output format and file path.
\item 用户：Use a ring of 12 modules, each 8 mm x 10 mm x 2 mm, radius 40 mm, clearance 1 mm, material water; point gamma source 0.662 MeV at (0,0,0) mm; physics QBBC; output json.
\item 系统：This would overwrite a geometry value you already confirmed.

**Pending Change:**
*   **Module Size:** From (7.83, 4.72, 6.86) mm to (8, 10, 2) mm.
*   **Ring Modules:** From 4 to 12.
*   **Ring Radius:** From 67.49 mm to 40 mm.
*   **Clearance:** From 0.19 mm to 1 mm.

Reply **'confirm'** to apply these changes, or **'keep existing'** to keep the current values.
\item 用户：confirm
\item 系统：To confirm the source direction, what are the x, y, and z components (dx, dy, dz) for the beam direction?
\end{itemize}
\textbf{内部链路（自然语言）：}
\begin{itemize}
\item 第 1 轮：动作=summarize\_progress，阶段=materials，闭环状态=False。
\item 缺失字段：materials.volume\_material\_map、source.type、source.particle、source.energy、source.position、source.direction、physics.physics\_list、output.format、output.path、geometry.ask.ring.module\_size、geometry.ask.ring.count、geometry.ask.ring.radius、geometry.ask.ring.clearance。
\item LLM 使用=True；回退=无；后端=llm\_semantic\_frame+runtime\_semantic。
\item 第 2 轮：动作=confirm\_overwrite，阶段=source，闭环状态=False。
\item 缺失字段：source.direction。
\item LLM 使用=True；回退=无；后端=llm\_slot\_frame+runtime\_semantic。
\item 第 3 轮：动作=ask\_clarification，阶段=source，闭环状态=False。
\item 缺失字段：source.direction。
\item LLM 使用=True；回退=无；后端=llm\_semantic\_frame+runtime\_semantic+confirmed\_pending\_overwrite。
\end{itemize}
\subsection{案例 MTM14（industrial\_panel）}
\textbf{对话：}
\begin{itemize}
\item 用户：我想做一个二维阵列探测板，但 pitch 和材料还没定。
\item 系统：目前我们已确定了使用二维阵列探测板作为几何结构类型，并完成了盒体尺寸的设置。

接下来还需要补充以下信息：
- **材料**：体积与材料绑定
- **源**：源类型、粒子类型、源能量、源位置
- **物理**：物理列表
- **输出**：输出格式、输出路径

请继续提供这些信息，以便完成配置。
\item 用户：?? 4 x 3 ?????????? 10 mm x 10 mm x 2 mm?pitch\_x 12 mm?pitch\_y 14 mm??? 1 mm??? ???? gamma 0.511 MeV??? (0,0,-80) mm??? +z????? QBBC??? root?
\item 系统：我检测到这会覆盖已经确认的值。待修改内容：几何结构类型：grid -> single\_box；盒体 x 方向尺寸：None -> 10.0。如果要应用这些修改，请直接回复“确认”；如果要保留原值，请直接说“保持原值”。
\item 用户：确认
\item 系统：配置已完成。当前已确认的参数包括：y 方向 pitch、x 方向 pitch、盒体 z 方向尺寸、盒体 y 方向尺寸。
\end{itemize}
\textbf{内部链路（自然语言）：}
\begin{itemize}
\item 第 1 轮：动作=summarize\_progress，阶段=materials，闭环状态=False。
\item 缺失字段：materials.volume\_material\_map、source.type、source.particle、source.energy、source.position、source.direction、physics.physics\_list、output.format、output.path、geometry.ask.grid.module\_size、geometry.ask.grid.count\_x、geometry.ask.grid.count\_y、geometry.ask.grid.pitch\_x、geometry.ask.grid.pitch\_y、geometry.ask.grid.clearance。
\item LLM 使用=True；回退=无；后端=llm\_semantic\_frame+runtime\_semantic。
\item 第 2 轮：动作=confirm\_overwrite，阶段=source，闭环状态=False。
\item 缺失字段：无。
\item LLM 使用=True；回退=无；后端=llm\_slot\_frame+runtime\_semantic。
\item 第 3 轮：动作=finalize，阶段=physics，闭环状态=True。
\item 缺失字段：无。
\item LLM 使用=True；回退=无；后端=llm\_semantic\_frame+runtime\_semantic+confirmed\_pending\_overwrite。
\end{itemize}
\subsection{案例 MTM15（shielding\_stack）}
\textbf{对话：}
\begin{itemize}
\item 用户：Build a stacked absorber sandwich along z; I will provide the layer sizes next.
\item 系统：Your geometry structure is now set to a stacked absorber sandwich along the z-axis. Here's what's still needed:

**Geometry:** Stack layer footprint and thicknesses.
**Materials:** Which volumes use which materials.
**Source:** Type, particle, and energy.
**Physics:** Physics list selection.
**Output:** Format and file path.
\item 用户：Use stacked layers with footprint 20 mm x 20 mm, thicknesses 2 mm, 4 mm, 6 mm, layer clearance 0.5 mm, outer box 30 mm x 30 mm x 20 mm, container clearance 1 mm, material aluminum; beam gamma 1.25 MeV from (0,0,-120) mm along +z; physics Shielding; output csv.
\item 系统：You are about to overwrite a previously confirmed geometry configuration.

**Pending Change:**
- **Stack footprint:** 16.33 mm x 10.28 mm → 20.0 mm x 20.0 mm
- **Layer thicknesses:** [6.68, 1.19, 5.59] mm → [2.0, 4.0, 6.0] mm
- **Layer clearance:** 0.73 mm → 0.5 mm
- **Outer box size:** 18.07 mm x 45.52 mm x 16.66 mm → 30.0 mm x 30.0 mm x 20.0 mm
- **Container clearance:** 0.87 mm → 1.0 mm

Reply **'confirm'** to apply this change, or **'keep existing'** to keep the current values.
\item 用户：confirm
\item 系统：Configuration complete. The following geometry items are confirmed: stack container clearance, stack container size, stack layer clearance, and stack layer thicknesses.
\end{itemize}
\textbf{内部链路（自然语言）：}
\begin{itemize}
\item 第 1 轮：动作=summarize\_progress，阶段=materials，闭环状态=False。
\item 缺失字段：materials.volume\_material\_map、source.type、source.particle、source.energy、source.position、source.direction、physics.physics\_list、output.format、output.path、geometry.ask.stack.footprint、geometry.ask.stack.thicknesses、geometry.ask.stack.layer\_clearance、geometry.ask.stack.parent\_size、geometry.ask.stack.nest\_clearance。
\item LLM 使用=True；回退=无；后端=llm\_semantic\_frame+runtime\_semantic。
\item 第 2 轮：动作=confirm\_overwrite，阶段=source，闭环状态=False。
\item 缺失字段：无。
\item LLM 使用=True；回退=无；后端=llm\_slot\_frame+runtime\_semantic。
\item 第 3 轮：动作=finalize，阶段=physics，闭环状态=True。
\item 缺失字段：无。
\item LLM 使用=True；回退=无；后端=llm\_semantic\_frame+runtime\_semantic+confirmed\_pending\_overwrite。
\end{itemize}
\subsection{案例 MTM16（nested\_target）}
\textbf{对话：}
\begin{itemize}
\item 用户：我想在一个盒子里嵌一个圆柱靶，先别急着补全。
\item 系统：好的，我们已确认了您要构建的“盒中嵌圆柱”几何结构。

目前，几何部分已明确结构类型，并确认了外层盒体的尺寸。接下来，我们还需要您提供：
1.  内嵌圆柱的尺寸（半径和半长）。
2.  为盒体和圆柱分别指定材料。
3.  设置粒子源的类型、种类和能量等信息。

请随时告诉我这些内容，我们继续配置。
\item 用户：外盒 120 mm x 120 mm x 120 mm，内嵌铜圆柱半径 20 mm、半长 40 mm，嵌套间隙 2 mm；材料 ?；束流 gamma 1 MeV，位置 (0,0,-150) mm，方向 +z；物理列表 QBBC；输出 json。
\item 系统：我注意到这将会覆盖已确认的几何参数。具体来说，外层盒子的尺寸、内嵌圆柱的半径与半长，以及嵌套间隙都将被更新为新值。

如果您确认应用这些修改，请直接回复“确认”；如果您希望保留原来的数值，请回复“保持原值”。
\item 用户：确认
\item 系统：好的，配置已完成。当前已确认的参数包括：嵌套间隙、内嵌圆柱半长、内嵌圆柱半径以及外层容器的尺寸。
\end{itemize}
\textbf{内部链路（自然语言）：}
\begin{itemize}
\item 第 1 轮：动作=summarize\_progress，阶段=materials，闭环状态=False。
\item 缺失字段：materials.volume\_material\_map、source.type、source.particle、source.energy、source.position、source.direction、physics.physics\_list、output.format、output.path、geometry.ask.nest.parent\_size、geometry.ask.nest.child\_radius、geometry.ask.nest.child\_half\_length、geometry.ask.nest.clearance。
\item LLM 使用=True；回退=无；后端=llm\_slot\_frame+runtime\_semantic。
\item 第 2 轮：动作=confirm\_overwrite，阶段=source，闭环状态=False。
\item 缺失字段：无。
\item LLM 使用=True；回退=无；后端=llm\_slot\_frame+runtime\_semantic。
\item 第 3 轮：动作=finalize，阶段=physics，闭环状态=True。
\item 缺失字段：无。
\item LLM 使用=True；回退=无；后端=llm\_semantic\_frame+runtime\_semantic+confirmed\_pending\_overwrite。
\end{itemize}
\subsection{案例 MTM17（coaxial\_shield）}
\textbf{对话：}
\begin{itemize}
\item 用户：I need a shell-style shielding toy model, but the shell thicknesses are still open.
\item 系统：Here's what we have so far:
*   **Geometry:** The structure is set to a shell. We still need the inner radius and the thicknesses for each shell layer.
*   **Materials:** We need to assign materials to each volume.
*   **Source:** We need the source type, particle type, energy, and position.
*   **Physics:** The physics list is configured.
*   **Output:** We need to choose an output format and file path.
\item 用户：Use a shell with inner radius 15 mm, thicknesses 5 mm and 8 mm, half length 40 mm, nested core radius 10 mm, nested core half length 30 mm, clearance 1 mm, material lead; isotropic neutron source 2 MeV at (0,0,0) mm; physics Shielding; output hdf5.
\item 系统：You're about to overwrite a geometry that was already set.

**Current geometry:**
*   **Shell:** Inner radius 7.15 mm, thicknesses [1.18, 3.43, 0.86] mm, half-length 22.36 mm.
*   **Nested Core:** Radius 11.60 mm, half-length 2.68 mm.
*   **Clearance:** 1.01 mm.

**New geometry:**
*   **Shell:** Inner radius 15.0 mm, thicknesses [5.0, 8.0, 0.83] mm, half-length 40.0 mm.
*   **Nested Core:** Radius 10.0 mm, half-length 30.0 mm.
*   **Clearance:** 1.0 mm.

Please reply **'confirm'** to apply this change, or **'keep existing'** to keep the current values.
\item 用户：confirm
\item 系统：Configuration complete. The following items are confirmed: shell clearance, nested core half length, nested core radius, and shell half length.
\end{itemize}
\textbf{内部链路（自然语言）：}
\begin{itemize}
\item 第 1 轮：动作=summarize\_progress，阶段=materials，闭环状态=False。
\item 缺失字段：materials.volume\_material\_map、source.type、source.particle、source.energy、source.position、source.direction、output.format、output.path、geometry.ask.shell.inner\_radius、geometry.ask.shell.thicknesses、geometry.ask.shell.half\_length、geometry.ask.shell.child\_radius、geometry.ask.shell.child\_half\_length、geometry.ask.shell.clearance。
\item LLM 使用=True；回退=无；后端=llm\_semantic\_frame+runtime\_semantic。
\item 第 2 轮：动作=confirm\_overwrite，阶段=source，闭环状态=False。
\item 缺失字段：无。
\item LLM 使用=True；回退=无；后端=llm\_slot\_frame+runtime\_semantic。
\item 第 3 轮：动作=finalize，阶段=physics，闭环状态=True。
\item 缺失字段：无。
\item LLM 使用=True；回退=无；后端=llm\_semantic\_frame+runtime\_semantic+confirmed\_pending\_overwrite。
\end{itemize}
\subsection{案例 MTM18（boolean\_cavity）}
\textbf{对话：}
\begin{itemize}
\item 用户：我想做一个挖空的盒体几何，但先只给你意图。
\item 系统：好的，我们已确认了您要创建“挖空的盒体”几何结构，并记录了盒体的尺寸。

接下来，为了完成这个布尔减运算的几何体，还需要您提供被挖去的那个形状（Solid B）在 x 和 y 方向上的尺寸。此外，整个几何体的材料也需要您来指定。

请随时告诉我这些信息。
\item 用户：???????A ? 40 mm x 40 mm x 40 mm ???B ? 20 mm x 20 mm x 20 mm ???? A ??? B??? ???? gamma 0.3 MeV??? (0,0,-60) mm??? +z????? QBBC??? xml?
\item 系统：我注意到，您这次提供的尺寸信息会让几何结构从“挖空盒体”（布尔减运算）变为“合并盒体”（布尔并运算），这会覆盖之前已确认的几何设置。

具体会修改：
- 几何类型：从“挖空盒体”改为“合并盒体”。
- 两个盒体的尺寸也会更新为您新提供的数值。

如果您确认要应用这些修改，请直接回复“确认”；如果您希望保留原来的“挖空盒体”几何，请回复“保持原值”。
\item 用户：确认
\item 系统：好的，我们已确认了“合并盒体”的几何尺寸、源项、物理列表和输出设置。

接下来，只需要为这个几何体中的各个体积指定材料，就可以生成完整的模拟程序了。
\end{itemize}
\textbf{内部链路（自然语言）：}
\begin{itemize}
\item 第 1 轮：动作=summarize\_progress，阶段=geometry，闭环状态=False。
\item 缺失字段：geometry.params.bool\_a\_x、geometry.params.bool\_a\_y、geometry.params.bool\_a\_z、geometry.params.bool\_b\_x、geometry.params.bool\_b\_y、geometry.params.bool\_b\_z、materials.volume\_material\_map、source.type、source.particle、source.energy、source.position、source.direction、physics.physics\_list、output.format、output.path、geometry.ask.boolean.solid\_a\_size、geometry.ask.boolean.solid\_b\_size。
\item LLM 使用=True；回退=无；后端=llm\_slot\_frame+runtime\_semantic。
\item 第 2 轮：动作=confirm\_overwrite，阶段=materials，闭环状态=False。
\item 缺失字段：materials.volume\_material\_map。
\item LLM 使用=True；回退=无；后端=llm\_slot\_frame+runtime\_semantic。
\item 第 3 轮：动作=summarize\_progress，阶段=materials，闭环状态=False。
\item 缺失字段：materials.volume\_material\_map。
\item LLM 使用=True；回退=无；后端=llm\_semantic\_frame+runtime\_semantic+confirmed\_pending\_overwrite。
\end{itemize}
\subsection{案例 MTM19（ring\_complete）}
\textbf{对话：}
\begin{itemize}
\item 用户：Create a complete ring setup: 16 modules, each 6 mm x 6 mm x 2 mm, radius 32 mm, clearance 0.5 mm, material silicon; point gamma source 0.511 MeV at (0,0,0) mm; physics QBBC; output json.
\item 系统：Here's your setup summary:

**Completed this turn:**
- **Geometry:** Structure, box X/Y sizes.
- **Materials:** Selected materials and their assignment to volumes.
- **Source:** Type, particle, and energy.
- **Physics:** Physics list.
- **Output:** Format and file path.

**Still needed:**
- **Source:** Direction.
\end{itemize}
\textbf{内部链路（自然语言）：}
\begin{itemize}
\item 第 1 轮：动作=summarize\_progress，阶段=materials，闭环状态=False。
\item 缺失字段：source.direction。
\item LLM 使用=True；回退=无；后端=llm\_slot\_frame+runtime\_semantic。
\end{itemize}
\subsection{案例 MTM20（grid\_complete）}
\textbf{对话：}
\begin{itemize}
\item 用户：???????????3 x 3 ????????? 12 mm x 12 mm x 3 mm?pitch\_x 15 mm?pitch\_y 15 mm??? 1 mm??? ?????? gamma 0.662 MeV??? (0,0,-90) mm??? +z????? QBBC??? root?
\item 系统：目前我们已经完成了几何结构、源设置、物理列表和输出配置的同步。接下来，只需要再补充一下体积与材料的绑定信息，就可以继续推进了。
\end{itemize}
\textbf{内部链路（自然语言）：}
\begin{itemize}
\item 第 1 轮：动作=summarize\_progress，阶段=materials，闭环状态=False。
\item 缺失字段：materials.volume\_material\_map。
\item LLM 使用=True；回退=无；后端=llm\_slot\_frame+runtime\_semantic。
\end{itemize}
\subsection{案例 MTM21（layered\_detector）}
\textbf{对话：}
\begin{itemize}
\item 用户：I want a stacked detector slab, but the individual layer thicknesses are not fixed yet.
\item 系统：You've set up the basic geometry structure for a stacked detector slab. Here's what we have and what's still needed:

**Completed:**
*   **Geometry:** The overall structure is configured as a stacked slab.

**Still Needed:**
*   **Geometry:** The footprint (size) and thicknesses for the individual stack layers.
*   **Materials:** Which material to assign to your detector volume.
*   **Source:** The type of source, the particle to simulate, and its energy.
*   **Physics:** The physics list to use.
*   **Output:** The format and file path for saving your results.
\item 用户：Use stack footprint 25 mm x 25 mm, thicknesses 1 mm, 2 mm, 1 mm, layer clearance 0.2 mm, outer box 30 mm x 30 mm x 12 mm, container clearance 0.5 mm, material silicon; beam electron 5 MeV from (0,0,-50) mm along +z; physics FTFP\_BERT; output json.
\item 系统：This would overwrite a geometry value you've already confirmed.

**Pending Change: Geometry**
*   **Stack Footprint:** 16.33 mm x 10.28 mm → 25.0 mm x 25.0 mm
*   **Stack Layer Thicknesses:** [6.68, 1.19, 5.59] mm → [1.0, 2.0, 1.0] mm
*   **Layer Clearance:** 0.73 mm → 0.2 mm
*   **Container Box Size:** 18.07 x 45.52 x 16.66 mm → 30.0 x 30.0 x 12.0 mm
*   **Container Clearance:** 0.87 mm → 0.5 mm

Please reply **'confirm'** to apply these new values, or **'keep existing'** to keep the current geometry.
\item 用户：confirm
\item 系统：Your configuration is complete. The following geometry items have been confirmed: stack container clearance, stack container size, stack layer clearance, and stack layer thicknesses.
\end{itemize}
\textbf{内部链路（自然语言）：}
\begin{itemize}
\item 第 1 轮：动作=summarize\_progress，阶段=materials，闭环状态=False。
\item 缺失字段：materials.volume\_material\_map、source.type、source.particle、source.energy、source.position、source.direction、physics.physics\_list、output.format、output.path、geometry.ask.stack.footprint、geometry.ask.stack.thicknesses、geometry.ask.stack.layer\_clearance、geometry.ask.stack.parent\_size、geometry.ask.stack.nest\_clearance。
\item LLM 使用=True；回退=无；后端=llm\_semantic\_frame+runtime\_semantic。
\item 第 2 轮：动作=confirm\_overwrite，阶段=source，闭环状态=False。
\item 缺失字段：无。
\item LLM 使用=True；回退=无；后端=llm\_slot\_frame+runtime\_semantic。
\item 第 3 轮：动作=finalize，阶段=physics，闭环状态=True。
\item 缺失字段：无。
\item LLM 使用=True；回退=无；后端=llm\_semantic\_frame+runtime\_semantic+confirmed\_pending\_overwrite。
\end{itemize}
\subsection{案例 MTM22（embedded\_sensor）}
\textbf{对话：}
\begin{itemize}
\item 用户：请创建完整的嵌套配置：外盒 80 mm x 80 mm x 80 mm，内嵌铅圆柱半径 15 mm、半长 25 mm，嵌套间隙 1 mm；材料 ?；点源 gamma 1 MeV，位置 (0,0,-100) mm，方向 +z；物理列表 Shielding；输出 csv。
\item 系统：配置已完成！当前已确认的设置包括：体积与材料绑定、输出路径、嵌套间隙（clearance）以及圆柱半径。
\end{itemize}
\textbf{内部链路（自然语言）：}
\begin{itemize}
\item 第 1 轮：动作=finalize，阶段=materials，闭环状态=True。
\item 缺失字段：无。
\item LLM 使用=True；回退=无；后端=llm\_slot\_frame+runtime\_semantic。
\end{itemize}
\subsection{案例 MTM23（layered\_shield\_complete）}
\textbf{对话：}
\begin{itemize}
\item 用户：Create a complete shell configuration: inner radius 12 mm, thicknesses 4 mm and 6 mm, half length 35 mm, nested core radius 8 mm, nested core half length 25 mm, clearance 1 mm, material iron; isotropic gamma source 0.8 MeV at (0,0,0) mm; physics QBBC; output hdf5.
\item 系统：Your configuration is complete. The following items have been confirmed: box dimensions (X, Y, Z) and the volume-to-material binding.
\end{itemize}
\textbf{内部链路（自然语言）：}
\begin{itemize}
\item 第 1 轮：动作=finalize，阶段=materials，闭环状态=True。
\item 缺失字段：无。
\item LLM 使用=True；回退=无；后端=llm\_slot\_frame+runtime\_semantic。
\end{itemize}
\subsection{案例 MTM24（boolean\_complete）}
\textbf{对话：}
\begin{itemize}
\item 用户：?????????????????A ? 30 mm x 30 mm x 20 mm ???B ? 10 mm x 10 mm x 40 mm ???? A ? union??? ???? gamma 0.5 MeV??? (0,0,-70) mm??? +z????? QBBC??? json?
\item 系统：目前我们已经完成了几何、源、物理和输出的基本设置，只剩下材料的“体积与材料绑定”需要补充。请告诉我每个几何体对应的材料，这样配置就完整了。
\end{itemize}
\textbf{内部链路（自然语言）：}
\begin{itemize}
\item 第 1 轮：动作=summarize\_progress，阶段=materials，闭环状态=False。
\item 缺失字段：materials.volume\_material\_map。
\item LLM 使用=True；回退=无；后端=llm\_slot\_frame+runtime\_semantic。
\end{itemize}
\end{document}